% Design justification for combined-dataset preprocessing choices.
% Intended to slot into the "Deviations from the Original Methodology"
% or a new "Combined-Dataset Preprocessing" subsection.

\paragraph{Naive concatenation.} Sequential concatenation of sessions with no
per-day correction, exactly matching Heelan et al.'s stated method of
sequentially concatenating recording sessions before splitting into
train/validation/test. Included as the baseline against which every other
variant is measured, and as the variant requiring no justification beyond
``this is what the original paper does.'' Fully causal: no statistic is
computed from data outside the calibration window used elsewhere in this
pipeline (Section~\ref{sec:deviations}).

\paragraph{Per-day normalization (whole-day statistics) -- rejected.} An
initial attempt standardized each day's features and targets independently
using that day's \emph{entire} timeline, including its own validation and
test bins, before concatenation. This variant is offline-only: a causal,
online system cannot compute a statistic from bins it has not yet observed.
It was tested purely to establish an upper bound on how much a per-day
correction could help, prior to committing to a deployable version. Once a
causal alternative (below) was tested and found to match this variant's
performance, this whole-day version was retired from the production
pipeline -- it added implementation complexity and a genuine online-decoding
liability without adding any measurable benefit over the causal version.

\paragraph{Per-day normalization (causal, calibration-window statistics) --
adopted.} The production variant computes each day's normalization
statistics (mean, standard deviation, and outlier-clipping bounds) from only
the leading 80\% of that day's timeline -- the same calibration-window
convention already used twice elsewhere in this pipeline: per-session
spike-detection threshold estimation (Section~\ref{sec:methods}), and the
global feature/target standardization Heelan et al.\ themselves describe
(``standardized to zero-mean using scikit-learn transformers fit on the
training data''). Those frozen statistics are then applied to that day's
full timeline, including its own validation and test bins, before
concatenation across days. This is directly analogous to a brief
calibration period at the start of a real recording session, after which
normalization is frozen and applied causally to all subsequent samples --
the standard pattern in deployed BCI decoders, and the pattern this pipeline
would need for online use regardless of preprocessing choice. Head-to-head
testing found this causal variant matched or marginally exceeded the
whole-day (offline-only) variant's validation and test correlation,
indicating the earlier gain reflects a real, calibratable per-day property
of the data rather than information leakage from future samples.

\paragraph{Outlier clipping.} Cross-day diagnostics (Section~\ref{sec:results})
identified a small number of extreme-magnitude neural feature bins
(0.1--1.1\% of values per day), concentrated on specific channels on
specific days and consistent with transient recording artifact rather than
broadly distributed task-related activity. Clipping each day's normalized
features to $\pm5$ standard deviations, using bounds derived from the same
calibration window described above, measurably improved validation and (more
substantially) test correlation over the unclipped causal variant. Because
the clip threshold is fixed and derived only from the calibration window,
this step is equally causal and deployable, and can be interpreted as a
standard artifact-rejection step rather than a form of post-hoc data
selection -- the threshold is not adapted per outlier, only calibrated once
per day from data available before the samples it is applied to.

\paragraph{Summary, in the context of a future online/SNN decoder.} Every
preprocessing choice retained in the production pipeline -- naive
concatenation and causal per-day normalization with clipping -- computes any
required statistic from a fixed calibration window that precedes the samples
it is applied to, and applies that frozen statistic identically to all
subsequent data. Neither choice requires knowledge of a session's future,
its validation/test split, or any other session's data. This was a
deliberate constraint: a preprocessing pipeline that only works because it
was allowed to peek at data an online system would not yet have is not a
useful foundation for an eventual real-time decoder, however well it scores
offline.
